
\chapter{Directional extensions and directional verbs} \label{ch:verbs}

%from Verkerk 2014:
%Wälchli	 (2009)	studies motion	 verb	
%choice	 in	 motion	 event	 descriptions	 in	 five	 path	 domains:	 enter,	 exit,	 ascend,	
%descend,	and	pass/cross.	He	uses	a	parallel	corpus	of	 translations	of the	Gospel	
%according	 to	Mark (a	 Bible	 text)	in	 a	 world-wide	 sample	 of	 117	languages.	 For	
%each	path	domain,	he	makes	an	assessment	of	how	often	path	verbs	are	used	in	
%the	 translations.	 Wälchli	 finds	 that	 about	 two-thirds	 of	 the	 languages	 in	 his	
%sample	use	path	verbs	to	encode	all	five	path	domains.

In this chapter, the expression of directional meaning in verbal extensions (as presented in the preceding chapters) is compared with the expression of directional meanings in verb roots. The comparison tests the hypothesis that, within the verbal system of a particular language, directional semantic content will be expressed either as a verbal extension or as a verb root, but not both. The results show that this holds true for some, but not all types of directional meaning within Chadic languages.

The pattern of complementary distribution is evident in the case of vertical directional meaning (upward and downward) and boundary-crossing directional meaning (inward and outward) as discussed in Sections \ref{sec:verticalV} and \ref{sec:boundedV}. In contrast, there is no evidence that ventive and itive directional extensions correlate with the absence of ventive or itive verb roots (or vice versa); however, there is a very strong one-way correlation with the presence of ventive extensions as a condition for the presence of a general motion verb with no directional meaning (Section \ref{sec:motionV}). Before discussing the results, Section \ref{sec:previous} reviews two previous publications comparing the expression of directional semantics in different parts of the grammar, and Section \ref{sec:motionverbs} provides an overview of how commonly various types of motion verbs are found in Chadic languages of each branch. Section \ref{sec:verbsummary} is a brief conclusion.

\section{Previous studies} \label{sec:previous}

The well-known studies of \citet{talmyLexicalizationPatternsSemantic1985,talmyCognitiveSemantics2000} on the syntax and semantics of motion have focused on the idea that directional information (path of motion) can be expressed either as part of the meaning of the main verb of a clause (verb-framed) or elsewhere in the clause (satellite-framed). The primary question in this line of research is about typical usage of whatever means are available for expressing information about the path of motion, and the data examined are corpora of extended speech acts (primarily narratives). 

In contrast, the research presented in this chapter is not concerned with usage, and therefore does not examine corpora. Rather, the focus is on what linguistic means are available for expressing specific types of directional meaning in the grammar of different languages. \citet[106--111]{verkerkEvolutionaryDynamicsMotion2014} reviews the relatively few previous publications which examine the means available for expressing path-related meaning across languages. Perhaps most prominently, \citet[458]{slobinMindCodeText1997} hypothesized that languages which characteristically express the path of motion in a part of the grammar other than the main verb (satellite-framed languages) will tend to have a higher number of manner-of-motion verbs in the lexicon. This proposal involves an underlying assumption that path and manner will rarely be co-lexicalized in the same verb root.\footnote{\citet[230]{slobinManyWaysSearch2004} discusses a Turkish verb glossed `climb' which is said to conflate manner-of-motion with upward direction. However, this is not the case for the English verb \textit{climb} \citep{levinLexicalizedMeaningManner2013}. The approach of  \citet{levinLexicalizedMeaningManner2013} treats manner/path complementarity as a subtype of manner/result complementarity.} The expression of path information in a satellite can be thought of as leaving available the main verb for expression of manner, which corresponds to having more manner-of-motion verbs in the lexicon \citep[463]{slobinMindCodeText1997}. This correlation suggests a dynamic relationship between what type of information is expressed in satellite forms, and what types of meanings are lexicalized in verb roots.

Whereas \citet{slobinMindCodeText1997} examines satellite-framed languages for trends in the lexicon, \citet[]{verkerkEvolutionaryDynamicsMotion2014} examines verb-framed languages hypothesizing that ``it could be the case that verb-framed languages have a larger lexicon of path of motion verbs (directional verbs) such as \textit{exit}, \textit{ascend}, and \textit{pass} as compared with satellite-framed languages'' \citep[111]{verkerkEvolutionaryDynamicsMotion2014}. This hypothesis is tested using a corpus of parallel translated texts for a set of 20 Indo-European languages to conclude that ``there exists a relationship between the motion encoding system... and the size of the manner verb lexicon and path verb lexicon'' \citep[137]{verkerkEvolutionaryDynamicsMotion2014}. 

The research question raised in this chapter is similar in that it is also about the distribution of directional meaning across different sub-components of the grammar of a particular language. The part of grammar under consideration is the relationship between lexicalization (verb roots) and grammaticalization (verbal extensions). As shown in the preceding chapters, most Chadic languages use a grammaticalized satellite form (verbal extensions) for the semantic function of expressing directional meaning. The hypothesis is that the existence of a verbal extension expressing a particular directional meaning correlates with the absence of a verb root expressing that same meaning in the lexicon (and vice versa). 

\section{Directional verbs} \label{sec:motionverbs}

While the research presented in this chapter investigates a similar issue to that explored by \citet{slobinMindCodeText1997} and \citet{verkerkEvolutionaryDynamicsMotion2014}, there are also some important differences. Rather than attempting to quantify an unbounded number of verb meanings, this study is restricted to six specific directional meanings.\footnote{The main reason for excluding manner-of-motion verbs is that the quality and quantity of lexical data varies considerably across the languages in the dataset. A study of manner-of-motion verbs would only include a small subset of the languages considered here, without any clear criteria for selection.} The meanings under investigation are the six directional meanings most commonly found in verbal extensions: ventive, itive, inward, outward, upward and downward (Chapters \ref{ch:identifying} and \ref{ch:distribution}). The restriction to these six directional meanings facilitates a granular comparison of the distribution of particular semantic concepts, rather than conflating these meanings in a count of the overall number of path-expressing motion verbs in various languages. 

Restricting the scope of examination to these six directional meanings allows for the inclusion of data from 90 Chadic languages. In comparison, \citet{verkerkEvolutionaryDynamicsMotion2014} uses a parallel translation corpus of 20 languages to quantify an unbounded number of motion verbs used by each language in the sample. This type of data is not immediately available for Chadic languages. One shortcoming of relying on lexicons and dictionaries (rather than texts) is that it is not possible to observe the usage of different linguistic means for expressing direction. This is particularly relevant in the cases where there are more than one means available for expressing the same directional meaning. For example, it is not completely redundant to have both a directional verb and a directional extension expressing the same path (compare English \textit{enter} and \textit{into}) as the two forms undoubtedly vary in their usage.

%It is also assumed, based on a few case studies, that the argument structure of motion verbs in Chadic languages can be distinguished from the argument structure of manner-of-motion verbs by how locative arguments are marked. For example, in Barayin, a manner-of-motion verb like \textit{d-ii} `walk' cannot take a locative argument, as shown in the ungrammatical phrase in (\ref{ex:bara01}). In contrast, directional motion verbs can take unmarked locative arguments, as in (\ref{ex:bara02}). In the case of Barayin, manner-of-motion and path are combined through the use of a serial verb construction, as in (\ref{ex:bara03}). (See \citet[251]{frajzyngierGrammarPero1989} for a similar observation about directional serial verb constructions in Pero.)

%\ea\label{ex:bara01} \langinfo{Barayin locative arguments}{}{\cite[141, 148]{lovestrandBarayinMorphosyntaxLexicalFunctional2022}} \\
%\ea 
%\gll \textbf{*} ane dee-gi (ŋ/iŋ) Andi \\
%{} {SBJ.1PL.EXCL} walk-{IPFV} ({PREP}/{ASOC}) Andi \\
%\glt `\textit{for:} We walked to Andi.'
%
%\ex \label{ex:bara02}
%\gll ka s-eyi Balal \\
%SBJ.\gl{3}\gl{sg}.M come-IPFV Balili \\
%\glt `He comes to/from Balili.'
%
%\ex \label{ex:bara03}
%\gll ane \textbf{dee-gi} \textbf{see-gi} Andi \\
%{} {SBJ.1PL.EXCL} walk-{IPFV} come-IPFV ({PREP}/{ASOC}) Andi \\
%\glt `We walked (here) to Andi.'
%\z
%\z

Three pairs of directional motion verbs are examined in this chapter: deictic motion verbs (ventive and itive), boundary-crossing motion verbs (inward and outward) and vertical motion verbs (upward and downward). In addition, this chapter also considers two related types of motion verbs. First, there are languages in which one of the basic motion verbs expresses general translocative motion without specifying any direction at all (as illustrated in Chapter \ref{ch:distribution}, Sections \ref{sec:vent1} and \ref{sec:itive1}). The same verb is used for ventive and itive motion. Second, there are verbs which express boundary-crossing motion but which are not specific as to whether the motion is inward or outward. The same verb can express either type of path.  

There are no clear-cut examples in the data from Chadic languages of a vertical motion verb which is underspecified as to whether the motion is upward or downward. A possible example of this type of verb is found in \hyperref[kera]{Kera} (East Chadic). The verb  \textit{lí} is glossed as `(\textit{be})\textit{steigen}, \textit{monter} [go up]'. However, in one example, shown in (\ref{ex:kera0}), this verb is translated as expressing downward motion, demonstrating that this verb apparently can be used for vertical motion in either direction. The ventive suffix in this example indicates that the motion is toward the deictic center, and evidently, in this context, the deictic center is understood to be at a lower elevation.\footnote{Note that \citet{ebertSpracheUndTradition1979} does not include interlinear glossing so these have been added to (\ref{ex:kera0}).} 

\ea\label{ex:kera0}
\langinfo{\hyperref[kera]{Kera} verb \textit{lí} `ascend/descend' with ventive suffix}{}{\cite[113]{ebertSpracheUndTradition1979}} \\
\gll lúŋ-dà	\\
ascend/descend-\gl{vent}	\\
\glt `descended' [German: `\textit{stieg herab}']
\z 

%data in tab:verbsbranch is from python code in file Verbs 
\begin{table}[t]
 \caption{Number of Chadic languages (by branch) with each type of motion verb}
	\label{tab:verbsbranch}
	\begin{tabularx}{\textwidth}{l Yrrr Y}
		\lsptoprule
{} &  West &  Biu-Mandara &  Masa &  East &  Total \\
\midrule
Languages in sample &    29 &           42 &     4 &    15 &     90 \\
\midrule
MOTION     &     5 &           12 &     2 &     0 &     19 \\
VENT    &    23 &           28 &     4 &    15 &     70 \\
ITIVE      &    25 &           30 &     3 &    14 &     72 \\
INTO/OUT     &     2 &            5 &     2 &     0 &      9 \\
INTO        &    25 &           16 &     2 &    15 &     58 \\
OUT        &    26 &           9 &     2 &    14 &     51 \\
UP         &    13 &           19 &     2 &    13 &     47 \\
DOWN       &    19 &           20 &     1 &    14 &     54 \\
		\lspbottomrule
	\end{tabularx}
\end{table}

\tabref{tab:verbsbranch} gives an overview of how many Chadic languages (by branch) have each type of motion verb attested in the available data. Data on motion verbs is available for 90 languages. The general motion verb (MOTION) and the boundary-crossing directional verb (INTO/OUT) are the least commonly found. Neither one is attested in East Chadic languages. Conversely, East Chadic languages are most likely to have the other six types of motion verbs. There is only one East Chadic language (of 15) without an itive verb, two without an upward verb and one without a downward verb (in the available data).

\section{Vertical directional verbs and extensions} \label{sec:verticalV}
\largerpage[2]
This section describes the cross-linguistic distribution of vertical directional meaning (upward and downward) whether expressed as part of the meaning of a motion verb root or as a verbal extension. There is a pattern of complementary distribution of upward and downward paths of motion in the verbal system of Chadic languages such that the vertical paths tend to be expressed either in the meaning of a verb or as a verbal extension, but not both. 

As discussed in Chapter \ref{ch:distribution}, Section \ref{sec:vertical2}, there are no vertical directional extensions outside of Biu-Mandara languages. Therefore, the other branches (West, Masa and East) will be excluded as it can be considered impossible for a vertical directional extension to co-occur with a vertical directional verb in those languages.\footnote{The number of vertical directional verbs in those branches is shown above in \tabref{tab:verbsbranch}. 13 of 15 East Chadic languages have an upward verb root. East Chadic languages generally do not have any directional extensions. Only 13 of 29 West Chadic languages have an upward verb root. West Chadic languages often have a ventive directional extension, but this does not correlate with the presence or absence of an upward verb root. A very similar pattern is found in regard to downward verb roots outside of Biu-Mandara languages.}
	
%data in tab:UPVandDIR is from python code in file VerticalVerbs 
\begin{table}
	\caption{Upward extensions in Biu-Mandara languages with and without an upward motion verb}
	\label{tab:UPVandDIR}
	\begin{tabular}{rrrc}
		\lsptoprule
		{} & UP ext & No UP ext & \% \\
		\midrule
		UP verb & 1 & 18   &  \bwpie{1}{18}\\
		No UP verb & 7 & 16  &  \bwpie{7}{16} \\
		\% &  \bwpie{1}{7} &  \bwpie{18}{16} & \\
		\lspbottomrule
	\end{tabular}
\end{table}

The matrix in \tabref{tab:UPVandDIR} shows how often upward directional verbs and upward directional extensions co-occur in Biu-Mandara languages.\footnote{The black and white pie charts in this table and similar tables throughout this chapter show the percentage of a present variable (e.g., has a UP verb) in black in relation to the total number of the relevant row or column the pie chart is aligned with.} There is just one language which has both an upward verb root and an upward extension, \hyperref[malg]{Malgwa}. In the other seven languages with an upward extension, there is no upward verb root. In Biu-Mandara languages without an upward directional extension, 18 have an upward verb root and 16 do not. Therefore, it is rare for any language to express upward directional meaning in both a verb root and a verbal extension.\footnote{Note that given the relatively low number of languages included in this subset of the data, it is not possible to test this pattern for statistical significance. However, see Section \ref{sec:verbsummary} for discussion of a statistically significant pattern when considering all vertical and boundary-crossing directional meanings.}
\clearpage

%data in tab:DOWNVandDIR is from python code in file VerticalVerbs 
\begin{table}
	\caption{Downward extensions in Biu-Mandara languages with and without a downward motion verb}
	\label{tab:DOWNVandDIR}
	\begin{tabular}{rrrc}
		\lsptoprule
		{} & DOWN ext & No DOWN ext & \% \\
		\midrule
		DOWN verb & 2 & 18 &  \bwpie{2}{18} \\
		No DOWN verb & 11 & 11 &  \bwpie{11}{11} \\
		\% &  \bwpie{2}{11} &  \bwpie{18}{11} &  \\
		\lspbottomrule
	\end{tabular}
\end{table}


A similar pattern can be seen in regard to downward motion, as shown in \tabref{tab:DOWNVandDIR}. Among Biu-Mandara languages with a downward verb root, two have a downward extension and 18 do not. In Biu-Mandara languages without a downward verb root, 11 have a downward extension and 11 do not. In conclusion, there is a pattern of near complementary distribution of vertical directional meaning in the verbal systems of Biu-Mandara languages, being expressed either as a verb stem or as a verbal extension.

The two languages with both a downward motion verb and a downward directional extension are \hyperref[lama]{Lamang} and \hyperref[gude]{Gude}. In the case of Gude, there are actually two verbal extensions that can express downward direction. The verbal extension \textit{-gərə} is shown to express either inward or downward direction. There is also an extension \textit{-paa} which ``is not productive in modern Gude and occurs only with a few words'' \citep[101]{hoskinsonGrammarDictionaryGude1983}.\footnote{\citet[359]{schuhChadicCornucopia2017} notes that historically the suffix ``derives from the locative adverb \textit{pà} ‘at the base of, down, under’ and a few of the verbs that bear this extension show a `downness' relation.'' }

One area of possible uncertainty in the data is the analysis of verbs that are glossed in English with the manner-of-motion verb \textit{climb},\footnote{See \citet{levinLexicalizedMeaningManner2013} for cases where the English verb \textit{climb} can be used either as a manner-of-motion verb or as a directional verb, but not as both simultaneously.} as in \hyperref[gude]{Gude} where the verb \textit{ndərə} `climb' can occur with an upward directional suffix, as in (\ref{ex:gude01}), or with a downward directional suffix, as in (\ref{ex:gude02}). Verbs glossed `climb' have not been included as vertical directional verbs, although the lexical semantics of these verbs in Chadic languages would ideally be subject to more scrutiny in future studies to make explicit whether they express vertical motion or manner-of-motion. 

\ea\label{ex:gude01}
\langinfo{\hyperref[gude]{Gude} verb \textit{ndərə} `climb' with upward suffix}{}{\cite[100]{hoskinsonGrammarDictionaryGude1983}} \\
\gll ndərə-gi	\\
climb-\gl{up}	\\
\glt `climb up'

\ex\label{ex:gude02}
\langinfo{\hyperref[gude]{Gude} verb \textit{ndərə} `climb' with downward suffix}{}{\cite[100]{hoskinsonGrammarDictionaryGude1983}} \\
\gll ndərə-gərə	\\
climb-\gl{down}/\gl{into}	\\
\glt `climb down'
\z 

\section{Boundary-crossing directional verbs and directional extensions} \label{sec:boundedV}

Turning to the distribution of inward and outward directional meaning, there is an inverse correlation (or near complementary distribution) between the presence of inward motion verbs and inward directional extensions, as well as between the presence of outward motion verbs and outward directional extensions. This is similar to the pattern seen in regard to vertical directional meaning (Section \ref{sec:verticalV}). Again, since boundary-crossing directional extensions are only found in Biu-Mandara languages, the other branches of Chadic languages will not be considered in this section.\footnote{The number of boundary-crossing directional verbs in those branches is shown above in \tabref{tab:verbsbranch}. Boundary-crossing verbs are ubiquitous outside of Biu-Mandara languages. All 15 East Chadic languages in the sample have both an inward verb root and an outward verb root. In West Chadic languages, 25 of 29 have an inward verb root and 26 of 29 have an outward verb root. Of the other three or four languages, two have a general boundary-crossing verb not specified for inward or outward direction.}

As shown in \tabref{tab:INTOVandDIR}, of the 16 languages with an inward directional extension, only three have an inward verb root. In the 26 languages without an inward directional extension, 13 have an inward verb root, and 13 do not. These numbers suggest a possible bias against the co-occurrence of inward verb roots and inward directional extensions in the same language.

%data in tab:INTOVandDIR is from python code in file BoundedVerbs 
\begin{table}
	\caption{Inward extensions in Biu-Mandara languages with and without an inward motion verb}
	\label{tab:INTOVandDIR}
	\begin{tabular}{rrrc}
		\lsptoprule
		{} & INTO ext & No INTO ext & \% \\
		\midrule
		INTO verb & 3 & 13  & \bwpie{3}{13} \\
		No INTO verb & 13 & 13   & \bwpie{13}{11} \\
		\% &  \bwpie{3}{13}	& \bwpie{13}{13} & \\
		\lspbottomrule
	\end{tabular}
\end{table}

The three languages that have both an inward verb root and an inward directional extension are \hyperref[budu]{Buduma}, \hyperref[dghw]{Dghwede} and \hyperref[marg]{Margi}. In Buduma, the inward suffix \textit{-al} is shown suffixed to an inward verb \textit{i} `\textit{hineingehen} [go in]', \textit{y-al}, with no change in meaning \citep[60]{lukasSpracheBudumaIm1939}. In the Dghwede data, there is just one example of the suffix \textit{-me} which appears to have inward directional meaning.

The case of \hyperref[marg]{Margi} is more complex. Margi has directional verbs which simultaneously express more than one dimension of the path of motion (cf. synexpression in Chapter \ref{ch:identifying2}, Section \ref{sec:portmanteau}). In the case of inward motion, the two verbs are the inward-itive \textit{lì} `go in, go (home)' and the inward-ventive \textit{sì} `come in, come (home)'. \citet[124, 149]{hoffmannGrammarMargiLanguage1963} notes the shared vowels in these verb forms, but treats this similarity as a question of etymology, not as part of the synchronic morphological system of Margi. As a further complexity, the inward directional extension in Margi only appears in contexts of transitive motion events (where the patient is the figure on a path of motion). In Margi, the inward verbs roots and the inward directional extension seem to have very different functions in the grammar. 

The patterns found in regard to outward direction are very similar to inward direction. Of 15 Biu-Mandara languages with an outward directional extension, only one also has an outward verb root, as shown in \tabref{tab:OUTVandDIR}. It is relatively common to find languages in which outward motion is not expressed in a motion verb or in a verbal extension (19 languages), but within the verbal system the pattern is one of complementary distribution where the outward directional meaning tends to be expressed either in a verb root or in a verbal extension, but not both.

%data in tab:OUTVandDIR is from python code in file BoundedVerbs 
\begin{table}
	\caption{Outward extensions in Biu-Mandara languages with an outward motion verb}
	\label{tab:OUTVandDIR}
	\begin{tabular}{rrrc}
		\lsptoprule
		{} & OUT ext & No OUT ext & \% \\
		\midrule
		OUT verb & 1 & 8  &  \bwpie{1}{8} \\
		No OUT verb & 14 & 19 &  \bwpie{14}{19} \\
		\% &  \bwpie{1}{14} &  \bwpie{8}{19} & \\
		\lspbottomrule
	\end{tabular}
\end{table}

The one language which has outward meaning expressed both as a verb root and also as a directional extension is \hyperref[marg]{Margi}, discussed above in regard to inward motion. There are two outward motion verbs in Margi: the outward-itive \textit{bà} `go out, go across' and the outward-ventive \textit{zə̀bà} `come out' \citep{hoffmannGrammarMargiLanguage1963}. The outward directional extension has the same form as the outward-itive verb root, which \citet[125]{hoffmannGrammarMargiLanguage1963} considers evidence of an etymological link. There is one example of the outward extension used with an intransitive motion verb, shown in (\ref{ex:marg6}). 

\ea\label{ex:marg6}
\langinfo{\hyperref[marg]{Margi} verb \textit{mə̀} `run' with the outward suffix \textit{-bá}}{}{\cite[122]{hoffmannGrammarMargiLanguage1963}} \\
\gll mə̀-bá	\\
run-\gl{out}		\\
\glt `to start up and run out'
\z 

In some Chadic languages, there is a motion verb that can be used to describe a path of motion across a boundary without lexically specifying whether the direction is inward or outward, as in (\ref{ex:gude100}). These verb roots will be called boundary-crossing motion verbs. Where the inward-outward distinction is not lexicalized in a verb root, it might logically be expected that this meaning would be expressed in a directional extension in order to provide a straightforward grammaticalized means of disambiguating whether a movement is inward (as in \ref{ex:gude101}) or outward (as in \ref{ex:gude102}). 

\ea\label{ex:gude1}
\ea\label{ex:gude100}
\langinfo{\hyperref[gude]{Gude} verb \textit{gim} `enter/exit' with ventive suffix }{}{\cite[100]{hoskinsonGrammarDictionaryGude1983}} \\
\gll gim-a	\\
enter/exit-\gl{vent}		\\
\glt `come in or out (toward speaker)'

\ex\label{ex:gude101}
\langinfo{\hyperref[gude]{Gude} verb \textit{gim} `enter/exit' with inward/downward suffix }{}{\cite[100]{hoskinsonGrammarDictionaryGude1983}} \\
\gll gim-a-gərə	\\
enter/exit-\gl{vent}-\gl{down}/\gl{into}		\\
\glt `come in (toward speaker)'

\ex\label{ex:gude102}
\langinfo{\hyperref[gude]{Gude} verb \textit{gim} `enter/exit' with outward/upward suffix }{}{\cite[100]{hoskinsonGrammarDictionaryGude1983}} \\
\gll gim-a-gi	\\
enter/exit-\gl{vent}-\gl{up}/\gl{out}		\\
\glt `come out (toward speaker)'
\z 
\z 

In the set of 90 Chadic languages, there are nine languages which have a boundary-crossing motion verb underspecified for inward or outward motion, as listed in \tabref{tab:inout}.\footnote{In addition, the East Chadic language \hyperref[kera]{Kera} uses the same verb stem to express inward and outward motion \citep{ebertSpracheUndTradition1976}. The verb root \textit{kélé} appears to always mean `enter' (inward motion) when used without a verbal extension. The verbal extension \textit{wə́ra} is described as a type of aspectual marker similar to ``totality'' markers described in many Chadic languages. To express outward motion, this verbal extension is combined with the verb \textit{kélé} `enter', in which case it becomes \textit{kél wə́ra} `exit'. However, since the verb \textit{kélé} on its own does not appear to be ambiguous between inward and outward motion meanings, it is not included in \tabref{tab:inout}.} Five of these languages are in the Biu-Mandara branch, two are Masa languages, and two are West Chadic languages. In languages with an underspecified boundary-crossing motion verb, there is not another verb root that expresses only inward motion or only outward motion.

%table tab:inout data from python script in file BoundedVerbs
\begin{table}
	\caption{Deictic and bound directional extensions in languages with a boundary-crossing motion verb}
	\label{tab:inout}
	\begin{tabular}{llllll}
		\lsptoprule
		Branch & Language 				&  VENT ext &  ITIVE ext &  INTO ext &  OUT ext \\
		\midrule
		Biu-Mandara &   \hyperref[musg]{Musgu}  &   Y &    Y &   N &  N \\
		Biu-Mandara &    \hyperref[buwa]{Buwal} &   Y &    Y &   N &  N \\
		Biu-Mandara &   \hyperref[mbud]{Mbudum} &   Y &    N &   N &  N \\
		Biu-Mandara &     \hyperref[gude]{Gude} &   Y &    N &   Y &  Y \\
		Biu-Mandara &     \hyperref[tera]{Tera} &   N &    N &   N &  N \\
		Masa &   		\hyperref[masa]{Masana} &   Y &    Y &   N &  N \\
		Masa &    \hyperref[muse]{Musey} 		&   Y &    Y &   N &  N \\
		West &	 \hyperref[shaa]{Ron-Scha} 		&   Y &    N &   N &  N \\
		West &    \hyperref[dass]{Dass} 		&   N &    N &   N &  N \\
		\lspbottomrule
	\end{tabular}
\end{table}

\tabref{tab:inout} indicates which of the four types of directional extensions are found in these nine languages. Contrary to what might be expected, only \hyperref[gude]{Gude} has an inward or outward directional extension (as illustrated in \ref{ex:gude1}). Ventive extensions are much more common. They are found in seven of these nine languages. It can be observed in several examples, such as (\ref{ex:buwa0}) and  (\ref{ex:mbud0}), that a ventive suffix may effectively disambiguate the inward-outward meaning assuming it is established whether the deictic center is inside or outside of the bounded space.

\ea\label{ex:buwa0}
\langinfo{\hyperref[buwa]{Buwal} verb \textit{dám} `enter/exit' with ventive suffix }{}{\cite[372]{viljoenGrammaticalDescriptionBuwal2013}} \\
\gll dám-xā á bzā	\\
enter/exit-\gl{vent} \gl{prep} outside		\\
\glt `Come outside. (Speaker is outside)'

\ex\label{ex:mbud0}
\langinfo{\hyperref[mbud]{Mbudum} verb \textit{təɗ} `enter/exit' with ventive suffix }{}{\cite[290] {fobakayyangPhonologieMorphologieSyntaxe2021}}  \\
\gll tə̀ɗ-àhà	\\
enter/exit-\gl{vent}	\\
\glt `Come out! [French: \textit{Sors !}]'
\z 

There are two languages, \hyperref[tera]{Tera} and \hyperref[dass]{Dass}, which have a boundary-crossing motion verb but do not have a ventive, inward or outward directional extension. Tera is reported to have an itive directional ``particle'' which is not considered a verbal extension, but which could conceivably be used in a manner parallel to the use of the ventive directional extension illustrated in (\ref{ex:buwa0}) and  (\ref{ex:mbud0}). There is limited information on these two languages, and so there are not any examples available which demonstrate how the inward-outward distinction might be made. It may be that in these languages the speakers use an adverbial phrase to make this distinction when necessary.

In conclusion, in regard to boundary-crossing motion, there is a tendency within the verbal system for inward and outward meaning to either be expressed in a verb root or to be expressed as a directional extension, but not both. However, in a minority of languages, boundary-crossing directional meaning is not expressed by either, and is presumably expressed adverbially when necessary. In languages that have an underspecified boundary-crossing motion verb which does not specify inward or outward direction, inward and outward directional extensions are surprisingly absent in all but one case. Where necessary, the inward-outward distinction could be made by the use of a ventive directional extension. 

\section{Ventive, itive and non-directional motion verbs}\label{sec:motionV}

Unlike vertical and boundary-crossing directional extensions, ventive directional extensions are found across all branches of Chadic languages (Chapter \ref{ch:distribution}, Section \ref{sec:vent2}). Ventive meaning is also very commonly expressed in a verb root. Whereas Sections \ref{sec:verticalV} and \ref{sec:boundedV} described an inverse correlation in the distribution of specific directional meanings expressed as verbal extensions or verb roots, in regard to ventive directional meaning, it is not uncommon for ventive directional extensions and ventive verb roots to co-occur in the same language. 

Of the 90 languages for which sufficient lexical data are available, 70 have a ventive verb root (\tabref{tab:verbsbranch}). As shown in \tabref{tab:VENTVandDIR}, among the 70 languages that have a ventive verb root, 29 also have a ventive directional extension. While there is still an overall tendency to avoid expressing ventive meaning in both forms, there is a sizeable minority of languages that do have both a ventive verbal extension and also a ventive verb root.\footnote{A Fisher Exact test shows the distribution to be statistically significant, p < 0.01.}

There are 17 languages with no ventive verb root. All of these languages instead have a motion verb which is not specified for any direction. As will be discussed in more detail below, all of these languages also have a ventive directional extension. This means that it is never the case that ventive meaning is absent from both verb roots and verbal extensions in any Chadic language (in this data set).

%data in tab:motionVandDIR is from python code in file Verbs 
\begin{table}
	\caption{Ventive extensions in Chadic languages with and without a ventive directional verb}
	\label{tab:VENTVandDIR}
	\begin{tabular}{rrrc}
		\lsptoprule
		{} & VENT ext & No VENT ext  & \% \\
		\midrule
		VENT verb 	 & 27 & 43  & \bwpie{27}{43} \\
		No VENT verb & 20 & 0  & \bwpie{20}{0} \\
		\%  & \bwpie{27}{20} & \bwpie{43}{0} &  \\
		\lspbottomrule
	\end{tabular}
\end{table}

\hspace*{-0.8pt}Unlike ventive directional extensions, itive directional extensions are restricted to Biu-Mandara and Masa languages with the exception of a marginal case in the East Chadic language \hyperref[kera]{Kera}. For the purpose of examining correlations between itive directional extensions and itive verb roots, only Biu-Mandara and Masa languages will be considered. As shown in \tabref{tab:ITIVEVandDIR}, of the 18 Biu-Mandara and Masa languages that have an itive directional extension 11 also have an itive verb root. In the 28 Biu-Mandara and Masa languages without an itive directional extension, 22 have an itive verb root. These numbers do not indicate any correlation between itive directional extensions and itive verb roots.

%data in tab:motionVandDIR is from python code in file Verbs 
\begin{table}
	\caption{Itive extensions in Biu-Mandara and Masa languages with and without an itive directional verb}
	\label{tab:ITIVEVandDIR}
	\begin{tabular}{rrrc}
		\lsptoprule
		{} & ITIVE ext & No ITIVE ext  & \% \\
		\midrule
		ITIVE verb 	 & 11 & 22  & \bwpie{11}{22} \\
		No ITIVE verb & 7 & 6  & \bwpie{7}{6} \\
		\%  & \bwpie{11}{7} & \bwpie{22}{6} &  \\
		\lspbottomrule
	\end{tabular}
\end{table}

As was mentioned in Chapter \ref{ch:identifying}, Sections \ref{sec:vent1} and \ref{sec:itive1}, a number of Chadic languages have a basic motion verb which is not lexically specified for any path of motion. The verb simply denotes the fact of motion, referred to here as MOTION. By default, the motion may be given an itive translation \citep[cf.][]{wilkinsWhenGOMeans1995}, as in (\ref{ex:mina3}), but the same verb is compatible with ventive semantics as well, such as when a ventive directional suffix is added, as in (\ref{ex:mina4}).

\ea\label{ex:mina3}
\langinfo{\hyperref[mina]{Mina} verb \textit{nd} `go/come' without ventive suffix}{}{\cite[174]{frajzyngierGrammarMina2005}} \\
\gll  à	ndə̀	zə́	vù \\
\gl{3}\gl{sg} go/come	\textsc{end-of-event}	\gl{q} \\
\glt `Will he go?'

\ex\label{ex:mina4}
\langinfo{\hyperref[mina]{Mina} verb \textit{nd} `go/come' with ventive suffix}{}{\cite[174]{frajzyngierGrammarMina2005}} \\
\gll  à	nd-á	zə́	vù \\
\gl{3}\gl{sg}	go/come-\gl{vent}	\textsc{end-of-event}	\gl{q} \\
\glt `Will he come?'
\z 

Motion verbs unspecified for direction are attested in 19 of the 90 languages for which information about the lexicon is available. These 19 languages are listed in \tabref{tab:motionV}. Most of the languages are of the Biu-Mandara branch, but two Masa languages and five West Chadic languages are also reported to have a general motion verb, as exemplified by \hyperref[pero]{Pero} in (\ref{ex:pero1}). All five of the West Chadic languages are of the West A2 group (Boleic and Tangalic languages).\footnote{In addition, in the East Chadic language \hyperref[lele]{Lele}, the verb \textit{èjè} `come' appears to be formed of the ventive extension \textit{jè} and the verb \textit{è} `go'. However, there is evidence that these forms have fused into a single verb root. So this is likely a language that formally had a MOTION verb which has lexicalized with the ventive extension as a ventive directional verb.}

\ea\label{ex:pero1}
\langinfo{\hyperref[pero]{Pero} verb \textit{wáatò} `go/come' with ventive suffix}{}{\cite[88]{frajzyngierGrammarPero1989}} \\
\gll  wát-tù	mìnáa-nò \\
go/come-\gl{vent}	house-\gl{1}\gl{sg} \\
\glt `come to my house'
\z

\begin{table}[p]
	\caption{Chadic languages with a general motion verb}
	\label{tab:motionV}
\begin{tabular}{lllll}
	\lsptoprule
	Branch & Language & MOTION V & ITIVE V & VENT V \\% &  UP & DOWN & OUT &  INTO& INTO/OUT  \\
	\midrule
	Biu-Mandara &   \hyperref[molo]{Moloko} &      Y &     N &       N  \\%&   Y &    Y &   Y &   Y &    \\
	Biu-Mandara &   \hyperref[gava]{Gavar} &      Y &     N &       N  \\%&   Y &    Y &   N &   N &  \\
	Biu-Mandara &   \hyperref[mbud]{Mbudum} &      Y &     N &       N  \\%&   N &    Y &   N &   N &      Y \\
	Biu-Mandara &   \hyperref[mina]{Mina} &      Y &     N &       N  \\%&   N &    Y &   N &   Y &    \\
	Biu-Mandara &   \hyperref[buwa]{Buwal} &      Y &     N &       N  \\%&   Y &    Y &   N &   N &      Y  \\
	Biu-Mandara &   \hyperref[cuvo]{Cuvok} &      Y &     N &       N  \\%&   Y &    Y &   Y &   Y &   \\
	Biu-Mandara &  	\hyperref[wuzl]{Wuzlam} &      Y &     N &       N  \\%&   Y &    N &   N &   N &   \\
	Biu-Mandara &   \hyperref[mbar]{Mbara} &      Y &     N &       N  \\%&   Y &    N &   Y &   Y &   \\
	Biu-Mandara &   \hyperref[musg]{Musgu} &      Y &   N &     N  \\%& N &  N & N & N &      Y \\
	Biu-Mandara &   \hyperref[mere]{Merey} &      Y &   N &     N  \\%& Y &  Y & N & N &      N \\
	Biu-Mandara	& 	\hyperref[muya]{Muyang}  &    Y &   N &     N  \\
	Biu-Mandara	& \hyperref[zulg]{Zulgo-Gemzek} 	&    Y &   N &     N \\
	Masa & 	\hyperref[muse]{Musey} & 		Y 	& Y	& Y \\
	Masa &   \hyperref[masa]{Masana} &   	Y 	&  N &       Y \\% &   N &    N &   N &   N &      Y \\
	West &   \hyperref[kare]{Karekare} &      Y &     Y &     N  \\%& N &    Y &   Y &   Y &    N  \\
	West &   \hyperref[tang]{Tangale} &      Y &     N &       N  \\%&   Y &    Y &   Y &   Y &    N \\
	West &   \hyperref[bole]{Bole} &      Y &     N &       N  \\%&   N &    Y &   Y &   Y &  \\
	West &   \hyperref[ngam]{Ngamo} &      Y &   N &     N  \\%&   Y &    Y &   Y &   Y &     \\
	West &   \hyperref[pero]{Pero} &      Y &   N &     N  \\%&   Y &    Y &   Y &   Y &  \\
	\lspbottomrule
\end{tabular}
\end{table}


%data in tab:motionVandDIR is from python code in file Verbs
\begin{table}[p]
	\caption{Ventive extensions in languages with and without a MOTION verb}
	\label{tab:motionVandVENT}
	\begin{tabular}{rrrc}
		\lsptoprule
		{} & VENT ext & No VENT ext  & \% \\
		\midrule
		MOTION verb 	 & 19 & 0  & \bwpie{19}{0} \\
		No MOTION verb & 28 & 43  & \bwpie{28}{43} \\
		\%  & \bwpie{19}{28} & \bwpie{0}{43} &  \\
		\lspbottomrule
	\end{tabular}
\end{table}


\tabref{tab:motionV} indicates whether each of these 19 languages also has an itive or ventive directional verb. Only three languages that have a general motion verb also have an itive or ventive directional verb. This includes the two Masa languages. \hyperref[masa]{Masana} has a ventive directional verb and \hyperref[muse]{Musey} has both itive and ventive directional verbs. The West Chadic language \hyperref[kare]{Karekare} has an itive directional verb.

Assuming a fundamental disposition to frame motion events in terms of deictic direction, it is expected that languages which do not have a ventive-itive directional distinction lexicalized as a verb stem will have some other means for making this distinction. As shown in \tabref{tab:motionVandVENT}, in Chadic languages, this distinction can always be made by means of a ventive directional extension.\footnote{A Fisher Exact test shows the distribution to be statistically significant, p < .01.} All 19 languages with a general motion verb root also have a ventive directional extension.

While the presence of a ventive directional extension is a necessary condition for a Chadic language to have a general motion verb, it is not the case that ventive directional extensions only occur where there is no ventive verb root (i.e., `come'). As seen in \tabref{tab:VENTVandDIR}, more than half of ventive directional extensions (27/47) are found in languages which also have a ventive directional verb.


In addition to all of the 19 languages with a general motion verb having a ventive directional extension, eight of them also have an itive directional extension. Taking into account that itive directional extensions are rarely found outside of Biu-Mandara and Masa languages, we can exclude West and East Chadic languages from consideration. As shown in \tabref{tab:motionVandTIVE}, among the 14 Biu-Mandara and Masa languages which have a general motion verb, eight of them have an itive directional extension (57\%). Among the 32 Biu-Mandara and Masa languages which do not have a general motion verb, ten have an itive extension (31\%). Therefore, there is no evidence of a correlation between the presence of a general motion verb and the presence of an itive directional extension.

%data in tab:motionVandDIR is from python code in file Verbs 
\begin{table}
	\caption{Itive extensions in Biu-Mandara and Masa languages with and without a MOTION verb}
	\label{tab:motionVandTIVE}
	\begin{tabular}{rrrc}
		\lsptoprule
		{} &  ITIVE ext & NO itive ext & \% \\
		\midrule
		MOTION verb 	 &  8 & 6 &  \bwpie{8}{6} \\
		No MOTION verb  & 10 & 22 &  \bwpie{10}{22} \\
		\% &  \bwpie{8}{10} &  \bwpie{6}{22} & \\
		\lspbottomrule
	\end{tabular}
\end{table}

In conclusion, among Chadic languages, general motion verbs are only found where there is also a ventive directional extension in the same language. However, the presence of a ventive directional extension does not necessarily correlate with the absence of a ventive directional verb root. Ventive direction is regularly found as both a verb and as an extension in the same language. There is no evidence of a correlation between the presence of itive directional extensions and general motion verbs. Likewise, the presence of an itive directional extension does not correlate with the absence of an itive directional verb, or vice versa. 

\section{Conclusion} \label{sec:verbsummary}

The hypothesis that directional meaning within the verbal system tends to be expressed either as a verb root or as a directional extension is found to be true for some types of directional meaning, but not all. The distribution of vertical and boundary-crossing directional paths confirms the hypothesis.\footnote{While it is not possible to check the statistical significance of these patterns for each specific semantic type of directional meaning, there is strong evidence of statistical significance when considering the relevant pattern of co-occurrence across all four types of vertical and boundary-crossing directional meanings. After conflating the data in Tables \ref{tab:UPVandDIR}, \ref{tab:DOWNVandDIR}, \ref{tab:INTOVandDIR} and \ref{tab:OUTVandDIR}, a chi-square test shows the distribution to be statistically significant: $\chi$² (1, N=160) = 21.3, p < .01 (without Yates' correction). The effect size is medium (Cohen's \textit{w} = 0.36) and statistical power is over 0.9.} These meanings tend to be either lexicalized as a verb root or as a verbal extension, but not both. 

This pattern can be viewed as the result of competition in the grammatical system between the use of a more specific directional verb root versus the use of a more general motion verb combined with a directional extension \citep{gardaniCompetitionMorphologyHistorical2019}. Once a vertical or boundary-crossing directional extension becomes part of the verbal morphology, the combination of that extension with the verb `go' is a possible synonym of a verb root that expresses an equivalent directional motion meaning.\footnote{The assumption is that vertical and boundary-crossing directional verb roots were previously present in those languages which later developed vertical and boundary crossing directional extensions. The reverse order, that directional extensions are more archaic, would be plausible if the forms of the relevant verb roots showed any sign of previously being morphologically composed before fossilizing into verb roots. One exception is \hyperref[bana]{Bana} but it unclear whether verbs forms are marked by nascent or erstwhile directional extensions.} Retention of the verb root would be expected if its use is frequent enough to block the use of the morphologically composed forms \citep{liebermanQuantifyingEvolutionaryDynamics2007,wuMorphologicalIrregularityCorrelates2019,smithRelationshipFrequencyIrregularity2023}. Evidently, this is not the case, and instead Biu-Mandara languages tend to drop the use of vertical or boundary-crossing directional verb roots in favor of regularized morphological patterns involving directional extensions.\footnote{No systematic studies of the relative frequencies of various verb roots across Chadic languages are available, but examination of a \hyperref[bara]{Barayin} corpus of about 30,000 words shows a large difference between the use of \textit{kolo} `go' (524) and \textit{sii}/\textit{ajo} `come' (368) compared with \textit{topo} `enter' (49), \textit{guso} `exit' (54), \textit{tado} `climb, ascend' (33) and \textit{jango} `descend' (17).}

In contrast with vertical and boundary-crossing, ventive and itive directional extensions are not found to be in complementary distribution with ventive and itive verb roots (respectively). Itive direction is the default interpretation of unmarked motion expressions \citep{wilkinsWhenGOMeans1995}, so itive directional extensions do not trigger the same competition as other types of direction in the verbal system. Ventive directional extensions can trigger the loss of a ventive directional verb root, but only under specific conditions. There are at least two reasons why the presence of a ventive directional extension does not usually trigger the loss of a ventive verb root. The first reason is that there may be semantic differences between verb roots that express ventive direction and verbal extensions. For example, in English, the ventive verb \textit{come} allows an interpretation in which the addressee is the deictic center: \textit{I'll come to you} \citep{goddardSemanticsComingGoing1997}. The same is not true of the adverbial \textit{here}. A similar semantic differences between ventive verb roots and ventive verbal extensions in Chadic languages could justify their co-occurrence.\footnote{This hypothesis implies that a similar semantic distinction would not be found when comparing verb roots and verbal extensions that express vertical or boundary-crossing directional meanings. Testing this hypothesis would require a significant amount of additional research on the lexical semantics of verb roots in Chadic languages.} The second reason is that there might not be a non-deictic verb root available for the ventive directional extension to combine with. In the case of vertical and boundary-crossing motion, these morphemes can freely combine with an itive verb root `go'. In contrast, a ventive verbal extension cannot combine with an itive verb root without a semantic clash. Only in languages which have a non-deictic verb root can there be a morphologically complex form to compete with the verb root `come'. In those cases where this does happen, the same process seems to take place -- the ventive verb root goes out of use, in spite of a relatively high rate of usage (See \tabref{tab:motionV} in Section \ref{sec:motionV}).
